\section{Is the Gap the Loss or the Data?}
\label{app:probe}

\S\ref{sec:discussion} claims the PTO/GRPO difference is about the data rather than the loss
family. This appendix gives the measurement behind that claim.

A note on what ``the data differ'' means here, since the candidate distributions are matched by
construction (\S\ref{sec:setup}): both methods draw $8$ completions from the same policy at the
same temperature $1.2$ at every branch point. What differs is the \emph{prefix} each group is
continued from --- an unmodified on-policy slice for GRPO, a trunk of compounding oracle-argmax
selections for PTO. The comparisons below vary whose recorded groups are used, so ``data'' in
the table means the states, not the sampler.

\subsection{Putting both updates on one probe}
Both methods take a group of candidate completions for a prompt, assign each a weight, and step
along the weighted sum. GRPO's weight is the standardized group advantage
$(r_g-\mu_g)/\sigma_g$; DPO's is $+1$ on the best candidate and $-1$ on the worst, zero
elsewhere. Rescaling to a common per-group size therefore puts the two on the same footing: for
each recorded group we compute the weighted-sum direction in a fixed sentence-embedding space and
compare directions by cosine.

Any such comparison is attenuated by how well each direction is estimated. We estimate an
\textbf{attenuation ceiling} by split-half --- the cosine between directions computed from two
disjoint halves of the same groups --- and report $\text{cosine}/\text{ceiling}$ as the corrected
agreement. A corrected value of $1.0$ means the two directions agree as well as the data permit.

\subsection{The decomposition}
The design is a $2\times2$: each method's own groups, weighted by each method's rule.

\begin{center}\small
\setlength{\tabcolsep}{3pt}
\begin{tabular}{@{}lrrr@{}}
\toprule
comparison & cos & ceil & corr. \\
\midrule
as trained (rule \emph{and} data differ) & 0.267 & 0.844 & \textbf{0.317} \\
\midrule
same data (PTO), rule swapped   & 0.908 & 0.793 & --- \\
same data (GRPO), rule swapped  & 0.988 & 0.929 & --- \\
\midrule
same rule (group-rel.), data differ & 0.356 & 0.896 & \textbf{0.397} \\
same rule (best/worst), data differ & 0.266 & 0.823 & \textbf{0.324} \\
\bottomrule
\end{tabular}
\end{center}

The reading is unambiguous. Holding the data fixed and swapping only the weighting rule leaves
the update direction essentially unchanged ($0.908$ and $0.988$ raw; these exceed their split-half
ceilings, which is what happens when two directions computed on the \emph{same} groups agree
better than two halves of those groups agree with each other, so we report them raw and do not
correct). Holding the rule fixed and varying whose groups are used leaves the methods as far apart
as they were as trained ($0.397$ and $0.324$, versus $0.317$). At matched look-ahead and a shared
oracle, DPO and group-relative weighting extract nearly the same direction from the same eight
completions.

\subsection{Selection pressure versus generated behaviour}
The same recorded groups let us separate what the update \emph{selects for} from what the policy
\emph{produces}. Per-iteration selection pressure toward affirmation --- the difference in
affirmation rate between what the update up-weights and what it down-weights --- runs
$-0.006\rightarrow+0.086\pm0.008$ (GRPO, iteration 10) and $0.008\rightarrow0.103\pm0.029$
(PTO, \emph{peaking at iteration 8}; PTO's iteration-10 value falls back to $0.039\pm0.038$ on
its thinnest group count). Over the same
iterations, the affirmation rate of what the policy generates moves $0.02\rightarrow0.54$ (GRPO)
and $0.04\rightarrow0.57$ (PTO), over-praise reaches $0.74$ of GRPO's candidates, and questions
per completion fall $0.71\rightarrow0.06$. The pressure is an order of magnitude smaller than the
drift it produces: it is compounding, not strong.

\subsection{Two caveats on this probe}
\paragraph{Per-iteration embedding directions are noisy.} An earlier version of this analysis
reported that a per-iteration direction predicts held-out preferences at $0.65\rightarrow0.71$.
That number was in-sample --- scored on the pairs it was fitted on. Held out, per-iteration PTO
directions win only $0.47$--$0.59$ with split-half reliability $0.15$--$0.32$. What survives is
the \emph{pooled} direction (split-half $0.597$ PTO, $0.911$ GRPO) and the exact lexical contrasts
above, which involve no embedding at all. The table in this appendix uses pooled directions
only.

\paragraph{PTO's training signal thins.} Branch points built per iteration fall $949\rightarrow410$
and the $\tau$-filter yield falls $0.82\rightarrow0.69$, so groups that actually contribute a
gradient fall $782\rightarrow281$, with the best-minus-worst margin decaying
$0.274\rightarrow0.196$. GRPO trains on $94$--$98\%$ of its groups throughout. A flattening PTO
curve may therefore be partly a data-starvation curve rather than a convergence curve --- which
matters for extrapolating beyond ten iterations, though not for the endpoint comparisons in the
body.

\subsection{Correlational link to the outcome}
With training iteration partialled out of both sides (the raw correlation is confounded with
iteration by construction), GRPO's change in MI-inconsistency tracks its affirmation push
($\rho=0.647$, $p=.043$), its length push ($0.706$, $p=.023$) and its over-praise push ($0.617$,
$p=.057$); PTO's does not ($-0.492$, n.s.). This is the same direction as the endpoint MICI gap
($0.84$ vs $0.49$), reached from the training data instead of the outcomes. With $n\leq10$
iterations per arm and no multiplicity correction, we report it as a mechanism consistent with the
curves, not as a cause.
